\section{Introduction}
\label{sec:intro}

A common pattern across production systems is a \emph{maintained sorted view}:
a collection of elements that is queried in sorted order (top-$k$, range scan,
or full ordered traversal) and updated in batches between queries. Limit order
books, polyphonic voice managers, packet schedulers, sprite renderers, and
radar trackers all share this shape.

The textbook choices are a self-balancing tree or a flat array re-sorted on
demand. \DeltaSort~\cite{deltasort} occupies a third point: it exploits the
observation that, in many of these systems, the set of \emph{updated indices is
already known} (the dirty list comes from the application layer for free), and
uses that information to re-sort the array in $O(n + k \log n)$ comparisons and
$O(n\sqrt{k})$ moves. For small $k$ relative to $n$, this beats both a full
sort and Extract--Sort--Merge (\ESM), the standard incremental-sort baseline.

A natural follow-up question is whether real systems actually live in the
small-$k$ regime, and if so, by how much. This paper answers that question
empirically.

\paragraph{Contributions.}
\begin{itemize}
  \item We port eight workloads spanning finance, audio, networking, defense,
        graphics, simulation, and pathfinding into a uniform Rust microbenchmark
        harness with rigorous statistics (500 samples, $p50/p95/p99$ with 95\%
        confidence intervals, two dirty-index distributions, verified
        correctness on every cell).
  \item We report that \DeltaSort wins 88\% of $(n,k)$ cells against \ESM,
        with peak speedups of 346$\times$ (game depth-sort) and 123$\times$
        (radar tracking) at workload-representative parameters. The 17 losses
        cluster at $k/n \geq 5\%$, matching the theoretical crossover.
  \item We validate the result in a real C++ codebase by porting \DeltaSort
        to a header-only template (fuzz-tested) and integrating with JUCE's
        \texttt{MidiMessageSequence}. \DeltaSort beats JUCE's stock
        \texttt{sort()} on all 40 measured cells, by up to
        \textbf{2624$\times$} for typical DAW edit operations.
  \item We identify one decisive negative: A$^*$ / Dijkstra / discrete-event
        simulation, where a binary heap with explicit position map beats
        \DeltaSort by 4--600$\times$. This bounds \DeltaSort's applicability
        domain and corrects a likely first instinct from the algorithm's
        complexity bound.
  \item We disclose a non-trivial tail-latency phenomenon: \DeltaSort's $p99/p50$
        ratio can reach 92$\times$ at specific $(n,k)$ configurations. Real-time
        systems must measure tails, not means, before adopting.
\end{itemize}
